% =============================================================================
% Chapter 11 — The ROM ($C000–$FFFF)
% =============================================================================
\chapter{The ROM}

\epigraph{``In ROM we trust --- for everything else, there's RAM.''}%
         {\textit{--- Anonymous systems programmer}}

\section*{Introduction}

The upper 16\,KB of the 6502 address space --- from \$C000 to \$FFFF ---
is occupied by read-only memory containing the EhBASIC interpreter and
monitor code.  Writes to this range are silently dropped, with two
important exceptions:

\begin{itemize}
  \item \textbf{SID registers} at \$D400--\$D43F are intercepted on write
        and routed to the SID sound chip emulation.  Reads still return ROM
        data.
  \item \textbf{Hardware vectors} at \$FFFA--\$FFFF are writable so that
        programs can install custom interrupt handlers and SID players can
        redirect execution.
\end{itemize}

The ROM image is loaded from \texttt{ehbasic.bin} at boot time.  It
provides the complete NovaBASIC environment: tokenizer, expression
evaluator, math library, string handling, I/O dispatch, and the extended
command set that distinguishes NovaBASIC from stock EhBASIC.


% ═══════════════════════════════════════════════════════════════════════════════
\section{ROM Layout Overview}
\label{sec:rom-layout}
% ═══════════════════════════════════════════════════════════════════════════════

\memrange{C000}{FFF9}{EhBASIC ROM}{R}{---}

The 16\,KB ROM contains the following regions (approximate boundaries ---
exact addresses depend on the build):

\begin{center}
\small
\begin{longtable}{@{}llp{8cm}@{}}
\toprule
\textbf{Range} & \textbf{Region} & \textbf{Contents} \\
\midrule
\endfirsthead
\toprule
\textbf{Range} & \textbf{Region} & \textbf{Contents} \\
\midrule
\endhead
\$C000--\$C002 & Cold start entry   & JMP to EhBASIC initialization \\
\$C003         & Warm start entry   & JMP to EhBASIC warm start \\
\$C000--\$DFFF & Interpreter core   & Tokenizer, executor, math routines,
                                       string handling, memory management \\
\$E000--\$EFFF & Functions \& I/O   & Built-in functions, commands, I/O
                                       handlers \\
\$F000--\$FF7F & Extended commands  & NovaBASIC additions (graphics, sound,
                                       networking, expansion memory, DMA,
                                       blitter, etc.) \\
\$FF80--\$FFF9 & Monitor code       & Power-on initialization and vector
                                       setup \\
\$FFFA--\$FFFF & Hardware vectors   & NMI, RESET, IRQ/BRK vectors
                                       (writable) \\
\bottomrule
\end{longtable}
\end{center}

\begin{notebox}
The boundaries above are approximate.  The interpreter, functions, and
extended commands are compiled from assembly source, and their exact
layout shifts with each build.  The entry points at \$C000, \$C003, and
\$FF80 are fixed by design.
\end{notebox}


% ═══════════════════════════════════════════════════════════════════════════════
\section{Key Entry Points}
\label{sec:rom-entry}
% ═══════════════════════════════════════════════════════════════════════════════

% --- $C000 LAB_COLD -----------------------------------------------------

\par\vspace{1.2em}%
\noindent
\begin{tikzpicture}[baseline=(addr.base)]
  \node[fill=retroblue, text=white, font=\ttfamily\bfseries\large,
        inner sep=4pt, minimum width=5.5em, anchor=west]
    (addr) {\$C000};
  \node[anchor=west, font=\sffamily\bfseries\large, text=retroink,
        xshift=6pt] at (addr.east)
    {LAB\_COLD};
  \node[anchor=east, font=\sffamily\small, text=retrogray]
    at ([xshift=\linewidth-1em]addr.west)
    {R\enspace|\enspace Default: ---};
\end{tikzpicture}%
\par\nopagebreak\vspace{0.3em}%
\label{addr:C000-cold}%
\index{\$C000@{\texttt{\$C000}} --- LAB\_COLD}%

EhBASIC cold start entry point.  Jumping here initializes all interpreter
state: clears variables, resets the program pointer, sets up the default
I/O vectors, prints the startup banner, and enters the command loop (the
\texttt{Ready} prompt).

This is the normal entry path after the monitor code completes its setup.
A cold start erases any BASIC program in memory.

% --- $C003 LAB_WARM ----------------------------------------------------

\par\vspace{1.2em}%
\noindent
\begin{tikzpicture}[baseline=(addr.base)]
  \node[fill=retroblue, text=white, font=\ttfamily\bfseries\large,
        inner sep=4pt, minimum width=5.5em, anchor=west]
    (addr) {\$C003};
  \node[anchor=west, font=\sffamily\bfseries\large, text=retroink,
        xshift=6pt] at (addr.east)
    {LAB\_WARM};
  \node[anchor=east, font=\sffamily\small, text=retrogray]
    at ([xshift=\linewidth-1em]addr.west)
    {R\enspace|\enspace Default: ---};
\end{tikzpicture}%
\par\nopagebreak\vspace{0.3em}%
\label{addr:C003}%
\index{\$C003@{\texttt{\$C003}} --- LAB\_WARM}%

EhBASIC warm start entry point.  Jumping here returns to the command
prompt \emph{without} erasing the current BASIC program in memory.
Variables are cleared, but program text is preserved.

This entry point is used by the \texttt{STOP} statement, the
\texttt{Break} key handler, and by user code that needs to abort execution
without losing the loaded program.


% ═══════════════════════════════════════════════════════════════════════════════
\section{Monitor Code (\$FF80--\$FFF9)}
\label{sec:monitor}
% ═══════════════════════════════════════════════════════════════════════════════

\memrange{FF80}{FFF9}{Monitor Code}{R}{---}

The reset vector at \$FFFC points to \$FF80.  This is the first code
executed at power-on or after a hardware reset.  The monitor performs the
following initialization sequence:

\begin{enumerate}
  \item \textbf{CLD} --- Clear decimal mode flag to ensure binary arithmetic.
  \item \textbf{LDX \#\$FF; TXS} --- Initialize the stack pointer to the
        top of the stack page (\$01FF).
  \item \textbf{Copy I/O vectors} --- Writes the default I/O vector table
        to page~2 (\$0200).  This table contains the addresses of the
        character input, character output, and load/save routines that
        EhBASIC calls through indirect jumps.
  \item \textbf{Check NOAUTO} --- Reads the \texttt{NOAUTO} environment
        variable.  If set, the monitor skips auto-run and proceeds
        directly to EhBASIC.
  \item \textbf{JMP LAB\_COLD} --- Transfers control to the EhBASIC cold
        start at \$C000.
\end{enumerate}

\begin{notebox}
The monitor code is deliberately minimal.  Its sole purpose is to bring
the machine to a known good state and hand off to the BASIC interpreter.
There is no built-in machine language monitor, hex editor, or disassembler
--- those functions are provided by the host debugger or the MCP server.
\end{notebox}


% ═══════════════════════════════════════════════════════════════════════════════
\section{SID Register Window}
\label{sec:rom-sid}
% ═══════════════════════════════════════════════════════════════════════════════

The SID sound chips are mapped at \$D400--\$D41F (SID~1) and
\$D420--\$D43F (SID~2) within this ROM range.  Although reads return ROM
data, writes are intercepted and routed to the SID emulation hardware.
See Chapter~12 for complete SID register documentation.

This dual-nature arrangement preserves compatibility with C64 SID player
files, which expect to write SID registers at \$D400, while allowing the
ROM image to occupy a contiguous 16\,KB block without gaps.

\begin{tipbox}
When reading back SID register values for diagnostic purposes, use the
SID shadow registers or the MCP server's \texttt{sid\_read} command.
Reading \$D400--\$D43F directly returns ROM bytes, not the values last
written to the SID chips.
\end{tipbox}


% ═══════════════════════════════════════════════════════════════════════════════
\section{Hardware Vectors (\$FFFA--\$FFFF)}
\label{sec:vectors}
% ═══════════════════════════════════════════════════════════════════════════════

\memrange{FFFA}{FFFF}{Hardware Vectors}{R/W}{Set at boot}

The six bytes at the very top of the address space contain the three
hardware interrupt vectors.  These are the \emph{only} writable bytes in
the entire ROM range.  They are writable so that SID players, custom
interrupt handlers, and timer-driven routines can redirect execution.

\begin{center}
\small
\begin{tabular}{@{}llp{7.5cm}@{}}
\toprule
\textbf{Address} & \textbf{Vector} & \textbf{Description} \\
\midrule
\$FFFA--\$FFFB & NMI   & Non-Maskable Interrupt handler address.  Triggered
                          by an external NMI signal; cannot be disabled by
                          the I flag. \\
\$FFFC--\$FFFD & RESET & Reset vector.  Points to \$FF80 (monitor code) at
                          boot.  The CPU reads this vector on power-on and
                          hardware reset. \\
\$FFFE--\$FFFF & IRQ/BRK & Maskable Interrupt / Break handler address.
                            Triggered by the IRQ line (when the I flag is
                            clear) or by a BRK instruction. \\
\bottomrule
\end{tabular}
\end{center}

Each vector is a 16-bit little-endian address.  The CPU reads these
vectors in hardware --- there is no software dispatch layer.

\begin{warningbox}
Overwriting the RESET vector will prevent the system from booting
correctly on next reset.  Only modify the NMI and IRQ vectors unless you
know what you are doing.
\end{warningbox}

\begin{lstlisting}
10 REM read the hardware vectors
20 PRINT "NMI:   $"; HEX$(DEEK($FFFA))
30 PRINT "RESET: $"; HEX$(DEEK($FFFC))
40 PRINT "IRQ:   $"; HEX$(DEEK($FFFE))
\end{lstlisting}

\begin{lstlisting}
10 REM install a custom IRQ handler at address $8000
20 REM (assumes handler code is already loaded there)
30 DOKE $FFFE, $8000
40 PRINT "IRQ VECTOR NOW POINTS TO $"; HEX$(DEEK($FFFE))
\end{lstlisting}

\begin{notebox}
The BASIC statements \texttt{IRQ} and \texttt{NMI} provide a safer,
higher-level way to install interrupt handlers.  They handle the vector
setup, register saving, and proper return sequence automatically.  Direct
vector manipulation is only needed for assembly-language interrupt
handlers.
\end{notebox}
